\chapter{La carte européenne}\label{ch:nw:the-european-map}

Le 6 juin 2022, Euronext a achevé le transfert de son centre de données principal, et de la \omterm{def:nw:colocation-products-and-how-they-are-sold:colo}{colocation} qui allait avec, de Basildon, à l'est de Londres, vers
le centre de données IT3 d'Aruba près de Bergame, dans le nord de l'Italie. Le bâtiment qu'il quittait appartenait à ICE~; celui vers lequel il a déménagé gère, selon les mots
du directeur général d'Euronext, un quart des volumes de négociation européens. Depuis Slough, où Euronext avait installé l'un de ses \omterm{def:nw:the-european-map:pop}{points
de présence} londoniens, la lumière dans le vide mettait \qty{258}{\micro\second} pour atteindre Basildon~; elle ne met désormais que
\qty{3.31}{\milli\second} pour atteindre Bergame. Les sociétés qui négocient les marchés d'Euronext depuis Londres ont vu s'ajouter, cette semaine-là, quatre millisecondes et demie
de fibre en ligne droite, et deux ans plus tard Euronext lui-même leur vendait une liaison hertzienne micro-ondes pour combler une partie de l'écart.

L'Europe n'a pas de New Jersey unique. Ses sites de négociation sont répartis entre Londres, Francfort, Zurich, le nord de l'Italie, Stockholm et Paris, et les
routes qui les relient (surtout Londres--Francfort) sont les courses à la latence du continent. Ce chapitre construit la moitié européenne du
tableau des sites avec les mêmes règles que le chapitre~10, ajoute le registre qui indique quelle plateforme apparie dans quel bâtiment à quelle date, et mesure ce que
la migration a produit.

\section{Les sites de la région londonienne}

\begin{definition}[Centre de données neutre]\label{def:nw:the-european-map:neutral}
Un \emph{centre de données neutre}\index{centre de données neutre} (carrier-neutral data centre) est un centre dont l'exploitant laisse n'importe quel opérateur de télécommunications amener son réseau
dans le bâtiment et vend des \omterm{def:nw:colocation-products-and-how-they-are-sold:mmr}{interconnexions} (cross-connects) entre n'importe quels locataires et n'importe quels opérateurs, plutôt que d'attacher les locataires à son propre réseau ou à un seul opérateur.
\end{definition}

Une plateforme peut se trouver dans un bâtiment neutre ou dans un bâtiment qui lui appartient~; ce qui compte pour une société est de pouvoir atteindre les systèmes de la plateforme et ses propres
opérateurs, depuis son armoire. Euronext décrit IT3 comme «~neutre par conception~», avec quatre points d'entrée des opérateurs (nord, sud, est et
ouest) et fibre noire vers le Milan Internet eXchange~; la \omterm{def:nw:colocation-products-and-how-they-are-sold:colo}{colocation} de Deutsche Börse à Francfort se trouve dans un bâtiment où les contrats d'hébergement
sont conclus entre la société et l'exploitant, Equinix, et la bourse y est «~l'un de nombreux clients~».

Londres compte trois sites pour ce chapitre. Le LD4 d'Equinix, sur le Slough Trading Estate à l'ouest de Londres, est l'endroit où Euronext a installé l'un de ses deux
\omterm{def:nw:the-european-map:pop}{points de présence} londoniens et où commence son service de liaison hertzienne micro-ondes vers Bergame~; il est aussi l'une des extrémités des routes Londres--Francfort ci-dessous. La
London Stock Exchange a transféré son centre de données principal, d'un site de Liverpool Street, dans la City, vers le campus de Telehouse à Docklands~: la bascule,
prévue en octobre 2022 pour Millennium Trading, Turquoise et le Group Ticker Plant, a été repoussée au 18 février 2023. Le
European Liquidity Centre d'ICE, à Basildon, dans l'Essex, fait tourner ICE Futures Europe et, jusqu'en juin 2022, accueillait aussi Euronext.

\begin{definition}[Point de présence]\label{def:nw:the-european-map:pop}
Un \emph{point de présence}\index{point de présence} (point of presence, PoP) d'un réseau est un lieu, en général des baies dans un \omterm{def:nw:the-european-map:neutral}{centre de données neutre}, où le
propriétaire du réseau (une plateforme, un opérateur) accepte les connexions des clients et achemine leur trafic vers ses propres sites, de sorte qu'un client
puisse atteindre une plateforme éloignée au moyen d'une \omterm{def:nw:colocation-products-and-how-they-are-sold:mmr}{interconnexion} dans un bâtiment voisin.
\end{definition}

Les trois sites londoniens sont proches~: 44~kilomètres de Slough à Docklands, 34 de Docklands à Basildon, 77 de Slough à Basildon, soit
147, 114 et \qty{258}{\micro\second} de lumière dans le vide. Dans un tel triangle, une société peut desservir les trois plateformes depuis
un seul bâtiment, avec quelques centaines de microsecondes de fibre vers chacune~; la migration a déporté l'un de ses angles vers l'Italie.

\begin{dated}{2026-09}{Qui apparie où en Europe}\label{dat:nw:the-european-map:sites}
\begin{itemize}
\item \textbf{Bergame} (Aruba Global Cloud Data Centre IT3, Via San Clemente 53, Ponte San Pietro)~: centre de données principal d'Euronext depuis le 6 juin
  2022, avec sa \omterm{def:nw:colocation-products-and-how-they-are-sold:colo}{colocation}, exploité par Euronext~; le site de reprise après sinistre reste dans la région parisienne~; \omterm{def:nw:the-european-map:pop}{points de présence} londoniens chez InterXion
  et Slough LD4.
\item \textbf{Francfort} (Equinix FR2, Kruppstrasse 121--127)~: les back-ends principaux d'Eurex et de Xetra et la \omterm{def:nw:colocation-products-and-how-they-are-sold:colo}{colocation} de Deutsche Börse.
\item \textbf{Zurich} (Equinix ZH4, Josefstrasse 225)~: la \omterm{def:nw:colocation-products-and-how-they-are-sold:colo}{colocation} de SIX Swiss Exchange, gérée par SIX~; un réseau de liaisons hertziennes micro-ondes de SIX relie ZH4
  à FR2.
\item \textbf{London Docklands} (campus de Telehouse, Coriander Avenue)~: le centre de données principal du London Stock Exchange Group depuis la
  migration depuis la City.
\item \textbf{Basildon} (European Liquidity Centre d'ICE)~: ICE Futures Europe~; l'adresse de la rue provient d'un annuaire et le site est approximatif.
\item \textbf{Upplands Väsby} (Smedbyvägen 6), à 25~kilomètres de Stockholm~: centre de données principal nordique de Nasdaq~; son site secondaire a été déplacé vers
  Equinix SK2, au sud de Stockholm, en 2025.
\end{itemize}
\end{dated}

\section{Francfort et Zurich}

La page de Deutsche Börse est précise~: le centre de données Equinix FR2 à Francfort est le centre de \omterm{def:nw:colocation-products-and-how-they-are-sold:colo}{colocation} «~dans lequel se trouvent les back-ends principaux d'Eurex et
de Xetra sont situés~». Ses connexions de \omterm{def:nw:colocation-products-and-how-they-are-sold:colo}{colocation} sont de 10~gigabits, avec les données de transaction et de marché sur des connexions distinctes afin que
elles ne fassent pas la queue l'une derrière l'autre (chapitre~1), tarifées au mois~: 6\,000~EUR pour une connexion Xetra de données de marché ou de saisie d'ordres, 8\,400
pour deux produits de données de marché ensemble, et 400~EUR pour son service de temps PTP. Les grilles tarifaires américaines du chapitre~9 sont déposées auprès de la SEC~; les
européennes sont publiées sous forme de tarifs, et elles sont tout aussi volatiles.

SIX exploite sa \omterm{def:nw:colocation-products-and-how-they-are-sold:colo}{colocation} dans l'Equinix ZH4 à Zurich et publie une latence moyenne de \qty{10.6}{\micro\second} pour une connexion en \omterm{def:nw:colocation-products-and-how-they-are-sold:colo}{colocation} de 10~gigabits
vers son moteur d'appariement (matching engine). Il exploite aussi un réseau de liaisons hertziennes micro-ondes de ZH4 vers FR2, à 10~mégabits par seconde, avec un accès égal pour tous les
participants~: une bourse qui vend le \omterm{def:nw:the-european-map:corridor}{couloir de latence} au bâtiment d'un marché voisin, comme le fait Euronext depuis Londres. Francfort et Zurich sont distantes
de \qty{306.7}{\kilo\metre}, \qty{1.02}{\milli\second} dans le vide et \qty{1.50}{\milli\second} en ligne droite dans la fibre.

\begin{table}[ht]
\centering\small
\begin{tabular}{llrrr}
\toprule
De & À & \omterm{def:nw:the-north-american-map:geodesic}{Distance géodésique} & Vide & Fibre \\
 & & (km) & (\si{\micro\second}) & (\si{\micro\second}) \\
\midrule
Slough LD4 & Francfort FR2 & 677.7 & 2\,260.6 & 3\,305.0 \\
Slough LD4 & Bergame IT3 & 992.2 & 3\,309.7 & 4\,838.7 \\
Francfort FR2 & Bergame IT3 & 497.6 & 1\,659.9 & 2\,426.8 \\
Slough LD4 & Basildon & 77.4 & 258.0 & 377.2 \\
Docklands & Basildon & 34.0 & 113.5 & 166.0 \\
Slough LD4 & Docklands & 44.0 & 146.9 & 214.8 \\
Basildon & Francfort FR2 & 603.2 & 2\,012.2 & 2\,941.8 \\
Basildon & Bergame IT3 & 936.2 & 3\,122.9 & 4\,565.7 \\
Francfort FR2 & Zurich ZH4 & 306.7 & 1\,023.1 & 1\,495.8 \\
Zurich ZH4 & Bergame IT3 & 204.5 & 682.2 & 997.4 \\
Francfort FR2 & Upplands Väsby & 1\,195.0 & 3\,986.0 & 5\,827.6 \\
Slough LD4 & Upplands Väsby & 1\,462.5 & 4\,878.2 & 7\,132.0 \\
\bottomrule
\end{tabular}
\caption{Distances et \omterm{def:nw:the-north-american-map:floor}{planchers de latence} à sens unique entre les sites européens de l'\cref{dat:nw:the-european-map:sites} (géodésiques WGS-84~; fibre
à $n_g = 1.4620$). Basildon, Bergame et Upplands Väsby sont placés au niveau de la rue, à un kilomètre ou deux des bâtiments, ce qui décale leurs
planchers d'une poignée de microsecondes. Données~: \texttt{nw\_europe.pair\_rows()}.}\label{tab:nw:the-european-map:pairs}
\end{table}

\begin{omfigure}
\begin{tikzpicture}
\begin{axis}[axis equal image,width=11.5cm,xmin=-560,xmax=1060,ymin=-440,ymax=380,xlabel={km à l'est de 4.5\textdegree{} E},
  ylabel={km au nord de 49\textdegree{} N},tick label style={font=\footnotesize,/pgf/number format/1000 sep={\,}},label style={font=\small},grid=major,grid style={gray!20},
  table/col sep=comma,unbounded coords=jump]
\addplot[omDef,thick,no markers] table[x=x,y=y]{figdata/networks/11-the-european-map/eu_links.csv};
\addplot[only marks,mark=*,mark size=2.4pt,omThm,visualization depends on={value \thisrow{anchor}\as\anc},nodes near coords,
  point meta=explicit symbolic,every node near coord/.append style={font=\footnotesize,anchor=\anc,black}]
  table[x=x,y=y,meta=name,col sep=comma,row sep=\\]{id,x,y,name,anchor\\ ld4,-375.101,279.029,Slough LD4,south\\ fr2,309.710,126.304,Francfort FR2 (Eurex~; Xetra),west\\ zh4,293.616,-178.256,Zurich ZH4 (SIX),east\\ bergamo,371.914,-364.528,Bergame IT3 (Euronext),west\\};
\addplot[only marks,mark=none,nodes near coords,point meta=explicit symbolic,
  every node near coord/.append style={font=\scriptsize,fill=white,inner sep=1pt,anchor=center,text=omMeth}]
  table[x=x,y=y,meta=text]{figdata/networks/11-the-european-map/eu_labels.csv};
\end{axis}
\end{tikzpicture}

\omcaption{Londres, Francfort, Zurich et Bergame à partir de coordonnées calculées, chaque liaison étiquetée par sa \omterm{def:nw:the-north-american-map:geodesic}{distance géodésique} et son temps de propagation dans
le vide~: le triangle Londres--Francfort--Bergame créé par la migration d'Euronext. Une projection équirectangulaire locale centrée sur
49\textdegree{}~N, 4.5\textdegree{}~E, destinée au seul dessin. Données~: \texttt{fig\_europe.py} sur \texttt{firm.geomap}.}\label{fig:nw:the-european-map:map}
\end{omfigure}

\section{La migration vers le nord de l'Italie}

Les questions fréquemment posées par Euronext sur le déménagement, telles qu'elles se présentaient en novembre 2021, décrivent l'ancien dispositif et le nouveau. Euronext
s'appuyait sur le centre de données de Basildon d'ICE, en vertu d'un contrat de service de longue durée~; la \omterm{def:nw:colocation-products-and-how-they-are-sold:colo}{colocation} était, selon le mot de son communiqué de presse, «~externalisée~».
Elle choisit un bâtiment tiers à Bergame, dans le pays de Borsa Italiana, qu'elle venait d'acheter, pour profiter de l'expérience de l'équipe
qui gérait la \omterm{def:nw:colocation-products-and-how-they-are-sold:colo}{colocation} de Borsa Italiana. Elle conserva son site de reprise après sinistre dans la région parisienne. Elle ouvrit deux \omterm{def:nw:the-european-map:pop}{points de présence} à Londres, chez
InterXion et à Slough LD4, depuis lesquels un client pouvait rejoindre Bergame, le site de reprise après sinistre et le centre de données de Borsa Italiana à Milan,
mais pas à Basildon. Elle ouvrit aussi le toit~: des fournisseurs sans fil et des sociétés de trading pouvaient installer des antennes sur une zone dédiée du toit, avec
Euronext garantissait une égalisation depuis l'arrière de l'antenne jusqu'à la \omterm{def:nw:colocation-products-and-how-they-are-sold:mmr}{salle d'interconnexion} (meet-me room), la règle du chapitre~9 appliquée à l'air.

\begin{omfigure}
\begin{tikzpicture}[font=\footnotesize,b/.style={draw,rounded corners=2pt,minimum height=0.6cm,align=center,inner sep=3pt},>=Stealth]
\node[b,fill=omDef!15] (fc) at (-0.8,0) {baie du client\\ à Slough LD4};
\node[b,fill=omThm!15] (pop) at (3.6,0) {Euronext\\ point de présence\\ (LD4)};
\node[b,fill=gray!10,minimum width=3.1cm] (net) at (7.6,0) {réseau Euronext\\ (fibre des opérateurs)};
\node[b,fill=omThm!25] (me) at (11.3,0.9) {Bergame IT3~:\\ moteurs d'appariement};
\node[b,fill=omThm!10] (dr) at (11.3,-0.9) {région parisienne~:\\ reprise après sinistre};
\draw[->,thick] (fc) -- node[above,font=\scriptsize]{interconnexion} (pop);
\draw[->,thick] (pop) -- (net);
\draw[->,thick] (net.east) -- (me.west);
\draw[->,thick] (net.east) -- (dr.west);
\node[b,fill=omMeth!12] (mw) at (3.6,1.9) {liaison hertzienne (EWIN)\\ antennes};
\draw[->,thick,omMeth,dashed] (fc.north) |- (mw.west);
\draw[->,thick,omMeth,dashed] (mw.east) -| node[pos=0.25,above,font=\scriptsize]{moins de 4 ms, secours par fibre} (me.north);
\node[font=\scriptsize,align=center,gray] at (7.6,-1.6) {un client d'un autre bâtiment rejoint le PoP par le circuit d'un opérateur};
\end{tikzpicture}

\omcaption{Comment une société londonienne rejoint Euronext après la migration, de façon schématique~: une \omterm{def:nw:colocation-products-and-how-they-are-sold:mmr}{interconnexion} en LD4 vers le \omterm{def:nw:the-european-map:pop}{point de présence} d'Euronext,
le réseau d'Euronext vers Bergame et vers le site de reprise après sinistre, ou, depuis 2024, le service hertzien d'Euronext depuis LD4 avec secours par fibre.
Sources~: la FAQ de la migration d'Euronext et son communiqué de presse de 2024 (l'\cref{dat:nw:the-european-map:sites}).}\label{fig:nw:the-european-map:pop}
\end{omfigure}

\begin{omfigure}
\begin{tikzpicture}
\begin{axis}[ybar,width=11.5cm,height=6.2cm,bar width=7pt,ymin=0,ymax=5.4,enlarge x limits=0.15,
  symbolic x coords={Slough LD4,Docklands,FR2,ZH4},xtick=data,ylabel={plancher aller simple vers Euronext (ms)},
  tick label style={font=\footnotesize},label style={font=\small},grid=major,grid style={gray!25},legend columns=4,
  legend style={font=\footnotesize,at={(0.5,-0.18)},anchor=north,draw=none,/tikz/every even column/.append style={column sep=6pt}},table/col sep=comma]
\addplot[fill=omDef!30,draw=omDef] table[x=from,y=vac_before]{figdata/networks/11-the-european-map/migration.csv};
\addplot[fill=omDef!80,draw=omDef] table[x=from,y=vac_after]{figdata/networks/11-the-european-map/migration.csv};
\addplot[fill=omThm!25,draw=omThm] table[x=from,y=fib_before]{figdata/networks/11-the-european-map/migration.csv};
\addplot[fill=omThm!75,draw=omThm] table[x=from,y=fib_after]{figdata/networks/11-the-european-map/migration.csv};
\legend{vide avant,vide après,fibre avant,fibre après}
\end{axis}
\end{tikzpicture}

\omcaption{Le \omterm{def:nw:the-north-american-map:floor}{plancher de latence} à sens unique depuis quatre sites vers le centre de données principal d'Euronext, avant (Basildon, jusqu'au 5 juin 2022) et après (Bergame,
à partir du 6 juin), dans le vide et dans une fibre rectiligne. Londres s'est éloignée de trois millisecondes dans le vide, de quatre millisecondes et demie dans la fibre~; Francfort
s'est un peu rapprochée et Zurich beaucoup plus. Données~: \texttt{nw\_europe.migration()} sur \texttt{firm.venuesites}.}\label{fig:nw:the-european-map:move}
\end{omfigure}

La \Cref{fig:nw:the-european-map:move} présente la migration en chiffres. Depuis Slough, le plancher dans le vide est passé de \qty{258}{\micro\second} à
\qty{3310}{\micro\second}, soit \qty{3.05}{\milli\second} de plus~; dans la fibre, de \qty{377}{\micro\second} à \qty{4839}{\micro\second},
soit \qty{4.46}{\milli\second} de plus. Depuis Docklands, c'est à peu près pareil. Depuis Francfort, le plancher a baissé de \qty{352}{\micro\second} (de 2.01 à
\qty{1.66}{\milli\second}), et depuis Zurich de \qty{1.81}{\milli\second}. La société londonienne qui avait négocié les marchés d'Euronext contre
ceux de Londres, distants de quelques centaines de microsecondes, les négocie désormais l'un contre l'autre depuis les deux extrémités d'un \omterm{def:nw:the-european-map:corridor}{couloir de latence} de trois millisecondes, et
choisit où placer les serveurs qui réagissent à chacun d'eux.

\omcode{code/networks/11-the-european-map/python/nw_europe.py}{53}{62}{Les planchers de la migration~: le site principal du registre pour Euronext Paris
la veille et le jour du déménagement, et les planchers vers chacun, calculés à partir d'une liste de sites.}\label{lst:nw:the-european-map:move}

\begin{proposition}[Déplacer le moteur déplace le triangle]\label{prop:nw:the-european-map:triangle}
Soit une société qui négocie sur deux plateformes situées aux sites $A$ et $B$ depuis un troisième site $F$. Le temps que met un signal pour aller d'une plateforme à l'autre en passant par la
société est au moins égal à $d(A,F) + d(F,B)$ divisé par la vitesse du milieu, et au moins égal à $d(A,B)$ divisée par cette même vitesse le long de la route directe~; la
position de la société ne dépend que de la première somme, qui est minimale lorsque $F$ se trouve sur la géodésique entre $A$ et $B$.
\end{proposition}

\begin{proof}
Inégalité triangulaire pour les géodésiques~: $d(A,F) + d(F,B) \ge d(A,B)$, avec égalité exactement lorsque $F$ se trouve sur le plus court chemin entre $A$ et
$B$.
\end{proof}

Avant la migration, Londres, Francfort et Euronext formaient une mince bande~: 678, 603 et 77~kilomètres, le moteur d'Euronext étant presque au-dessus de
Londres. Après celle-ci, le triangle est large~: 678, 498 et 992~kilomètres. Une société qui relaie entre Francfort et Bergame en passant par Londres paie
$678 + 992 = 1\,670$ kilomètres de trajet pour une distance directe de 498~; une société qui veut se placer entre Euronext et Eurex devrait se placer sur la
ligne Francfort--Bergame, là où Zurich se trouve presque (307 plus 205~kilomètres, 511 contre 498).

\section{Les pays nordiques et le reste}

Le centre de données principal nordique de Nasdaq se trouve à Upplands Väsby, à 25~kilomètres au nord de Stockholm, et son site secondaire a été déplacé en 2025 chez Equinix
SK2, au sud de la ville. Il est à \qty{1462.5}{\kilo\metre} de Slough (\qty{4.88}{\milli\second} dans le vide) et à \qty{1195}{\kilo\metre} de
Francfort~: le site le plus éloigné de Londres parmi les sites de ce chapitre, et un site auquel reviennent les chapitres~22 et~23.

Le reste de l'Europe suit le même schéma et la même méthode~: les documents de chaque plateforme désignent un bâtiment, souvent dans un \omterm{def:nw:the-european-map:neutral}{centre de données
neutre} (carrier-neutral) de sa propre ville~; le registre contient une ligne par plateforme, par rôle et par période, avec sa source. Le communiqué d'Euronext de 2024 évoque une seule
pool de liquidité sur une plateforme de négociation unique couvrant sept marchés européens, avec son centre de données central à Bergame~; le registre y consigne
les six marchés que nomment les sources de la migration, et laisse de côté celui de Borsa Italiana tant qu'une source ne date pas son transfert.

\omcode{code/firm/venuesites/firm_venuesites.py}{78}{98}{Les requêtes du registre plateforme-site~: lignes par plateforme, par rôle et par date~; le site principal d'une plateforme à une
date donnée (une erreur si la table n'en donne aucun ou plusieurs)~; les plateformes correspondant à un même bâtiment~; et le site de chaque plateforme classé par distance depuis
un bâtiment donné.}\label{lst:nw:the-european-map:registry}

\begin{omfigure}
\begin{tikzpicture}
\begin{axis}[xbar,width=11cm,height=6cm,bar width=6pt,xmin=0,xmax=7.8,ytick=data,
  yticklabels={Docklands,Basildon,FR2,ZH4,Bergame,Vasby},y dir=reverse,
  xlabel={plancher de latence à sens unique depuis Slough LD4 (ms)},tick label style={font=\footnotesize},label style={font=\small},
  grid=major,grid style={gray!25},legend pos=north east,legend style={font=\footnotesize},table/col sep=comma]
\addplot[fill=omDef!55,draw=omDef] table[y expr=\coordindex,x=vacuum_ms]{figdata/networks/11-the-european-map/from_ld4.csv};
\addplot[fill=omThm!40,draw=omThm] table[y expr=\coordindex,x=fibre_ms]{figdata/networks/11-the-european-map/from_ld4.csv};
\legend{vide,fibre}
\end{axis}
\end{tikzpicture}

\omcaption{Chaque site de plateforme européen du registre, classé selon son \omterm{def:nw:the-north-american-map:floor}{plancher de latence} à sens unique depuis Slough LD4, dans le vide et en fibre rectiligne~: les sites londoniens
à moins d'un tiers de milliseconde, Francfort et Zurich entre deux et quatre, Bergame et Upplands Väsby au-delà. Données~: \texttt{fig\_europe.py},
\texttt{Registry.nearest} depuis LD4.}\label{fig:nw:the-european-map:ld4}
\end{omfigure}

\section{Le couloir Londres--Francfort}

\begin{definition}[Couloir de latence]\label{def:nw:the-european-map:corridor}
Un \emph{couloir de latence}\index{couloir de latence} (latency corridor) est une paire de sites de négociation entre lesquels plusieurs réseaux se disputent le transport de flux sensibles à la latence,
chacun vendant son tracé selon sa latence à sens unique~: Londres--Francfort, Chicago--New Jersey, Zurich--Francfort.
\end{definition}

Le plancher du couloir est fixé par la géodésique~: \qty{677.7}{\kilo\metre} de Slough LD4 à Francfort FR2, \qty{2.261}{\milli\second} dans l'air et
\qty{3.305}{\milli\second} en fibre rectiligne. En mars 2015, McKay Brothers a annoncé un tracé entre les deux bâtiments à
\qty{4.64}{\milli\second} aller-retour, \qty{2.32}{\milli\second} à sens unique~: un \omterm{def:nw:the-north-american-map:factor}{facteur de tracé} de 1.026 dans l'air, \qty{59}{\micro\second} au-dessus du
plancher, sur un trajet qui doit traverser la Manche. Il était déjà plus rapide d'une milliseconde que ne pourrait jamais l'être la fibre rectiligne.

\begin{table}[ht]
\centering\small
\begin{tabular}{lrrrrr}
\toprule
Tracé hertzien & Géodésique & Publié & Plancher & Facteur & Plancher fibre \\
 & (km) & (ms) & dans l'air (ms) & de tracé & (ms) \\
\midrule
LD4 -- FR2 (2015) & 677.7 & 2.32 & 2.261 & 1.026 & 3.305 \\
LD4 -- Bergame IT3 (2024) & 992.2 & $< 4$ & 3.311 & $< 1.21$ & 4.839 \\
\bottomrule
\end{tabular}
\caption{Tracés hertziens européens publiés comparés à leurs planchers~: l'aller-retour 2015 de McKay Brothers divisé par deux, et le chiffre 2024 d'Euronext pour
EWIN, qui est une borne supérieure, de sorte que son \omterm{def:nw:the-north-american-map:factor}{facteur de tracé} l'est aussi. Données~: \texttt{nw\_europe.route\_rows()}.}\label{tab:nw:the-european-map:routes}
\end{table}

Le couloir Londres--Bergame est celui d'Euronext. Son service EWIN, lancé en juillet 2024 avec McKay Brothers, achemine des ordres par micro-ondes depuis
LD4 vers IT3 en «~moins de 4 millisecondes~», avec une sauvegarde intégrale en fibre. Le plancher dans l'air est de \qty{3.311}{\milli\second}, de sorte que la borne publiée
autorise un \omterm{def:nw:the-north-american-map:factor}{facteur de tracé} pouvant atteindre 1.21~; la fibre rectiligne ne pourrait pas faire mieux que \qty{4.839}{\milli\second}. Une borne publiée n'est pas une
mesure~: elle indique à une société ce que la plateforme garantit, et la société mesure le reste (chapitre~15).

\begin{remark}[Ordres dans un sens, données dans l'autre]\label{rem:nw:the-european-map:direction}
EWIN achemine des ordres de Londres à Bergame~; la société a toujours besoin des données de marché d'Euronext à Londres, par fibre depuis les \omterm{def:nw:the-european-map:pop}{points de présence} ou par
toute autre voie qu'elle achète. Un couloir est constitué de deux courses à la latence, une par sens, et une société peut les acheter à des réseaux différents.
\end{remark}

\begin{method}[Ajouter une plateforme au registre]\label{met:nw:the-european-map:add}
\begin{enumerate}
\item Repérer le MIC de la plateforme dans la liste ISO 10383, et son site dans les documents de connectivité propres à la plateforme ou dans ses dépôts.
\item Ajouter le bâtiment à la table des sites avec sa source d'adresse et son géocodage (méthode du chapitre~10), et
  si le géocodage ne vaut qu'au niveau de la rue, une note de précision.
\item Ajouter une ligne au registre par rôle (principal, secondaire, reprise après sinistre, \omterm{def:nw:the-european-map:pop}{point de présence}, diffusion des données de marché), avec sa source
  et la date de consultation~; une migration compte deux lignes, l'ancienne fermée la veille de l'ouverture de la nouvelle.
\item Recalculer les distances et les planchers, et dater l'ensemble de la table~: les plateformes changent de site, comme l'ont fait Euronext et le London Stock Exchange en moins d'un an.
\end{enumerate}
\end{method}

\omcode{code/firm/venuesites/firm_venuesites.py}{51}{55}{Validité d'une ligne du registre à une date~: un début ouvert ne dit rien des dates antérieures~;
une fin ouverte signifie que la ligne est toujours valide.}\label{lst:nw:the-european-map:dates}

\section{Tutoriel~: le registre européen}\label{tut:nw:the-european-map:tut}

\begin{tutorial}
\textbf{Objectif.} Étendre la table des sites à l'Europe, construire le registre plateforme-site, et mesurer la migration d'Euronext.
\textbf{État final~:} \cref{fig:nw:the-european-map:map,fig:nw:the-european-map:move,fig:nw:the-european-map:ld4} et
\cref{tab:nw:the-european-map:pairs,tab:nw:the-european-map:routes}.
\begin{tutsteps}
\item \textbf{Sites.} Sept lignes européennes dans \texttt{data/networks/sites.csv}, chacune avec ses entrées au grand livre~; trois sont au niveau de la rue et le signalent.
\item \textbf{Registre.} \texttt{code/firm/venuesites/data/venue\_sites.csv}~: une ligne par plateforme, rôle et période~;
  \texttt{Registry.load()} refuse une ligne sans source, avec un site inconnu, ou dont les dates sont inversées.
\item \textbf{La migration.} \texttt{nw\_europe.migration()} (\cref{lst:nw:the-european-map:move}) demande au registre le site principal d'Euronext
  le 5 et le 6 juin 2022, puis calcule les planchers~; \texttt{Registry.moves} liste le déplacement lui-même pour XPAR.
\item \textbf{Jonction avec le livre~1.} \texttt{join\_venues} apparie le registre à l'échantillon \texttt{firm.venues} du livre~1 par MIC, sans
  le modifier~; les MIC absents de l'échantillon sont ignorés, non inventés.
\item \textbf{Figures.} \texttt{python fig\_europe.py} écrit la carte, les barres de la migration et les distances depuis LD4.
\end{tutsteps}
\textbf{Pour aller plus loin.} Ajouter le centre de données de Milan de Borsa Italiana et le site de reprise après sinistre de la région parisienne une fois leurs bâtiments
documentés~; interroger \texttt{Registry.nearest} depuis Francfort plutôt que depuis Slough.
\end{tutorial}

\section{Construction~: le registre plateforme-site}\label{bld:nw:the-european-map:venuesites}

\begin{build}
\bfield{Rôle} Quelle plateforme correspond, quel site de reprise après sinistre elle conserve, et dans quel bâtiment, à quelle date et avec quelles sources elle accepte des connexions~:
l'entrée des chapitres~22 à~26 sur la connectivité et du plan du chapitre~29.
\bfield{Interface} \texttt{firm\_venuesites}~: \texttt{VenueSite(mic, venue, role, site, valid\_from, valid\_to, source, as\_of)},
\texttt{Registry.load}, \texttt{rows(mic, role, on)}, \texttt{primary(mic, on)}, \texttt{venues\_at(site, on)},
\texttt{nearest(site, role, on, medium)}, \texttt{moves(mic)}, \texttt{stale(today, days)}, \texttt{join\_venues(registry, venues)}.
\bfield{Règles} Chaque ligne a une source et une date de consultation~; chaque site figure dans la table de \texttt{firm.geomap}~; une plateforme a exactement un site principal à
toute date, sinon la requête échoue~; le registre du livre~1 est lu, jamais modifié.
\bfield{Tests d'acceptation} \texttt{code/firm/venuesites/tests/}~: la migration d'Euronext par date, les plateformes les plus proches de LD4 dans l'ordre, la
jonction avec l'échantillon du livre~1, et chaque ligne mal formée refusée.
\bfield{Pour aller plus loin} Lire la table des sites et le registre depuis un fichier unique versionné, avec un historique des modifications~; ajouter pour chaque plateforme ses points de diffusion des données de marché
et leurs \omterm{def:nw:networking-for-trading:group}{groupes multicast}.
\end{build}

\begin{omsources}
\item Euronext, communiqués des 15 juin 2022 (Core Data Centre migration) et 11 juillet 2024 (EWIN)~; \emph{Frequently Asked Questions about the
Euronext Data Centre Migration}, v6.1 (novembre 2021).
\item Deutsche Börse, Xetra co-location services~; SIX, co-location service fees (janvier 2026) et page de connectivité~; Nasdaq, Primary
Datacenter Stockholm (Väsby) et l'avis 2025 sur le site secondaire~; ICE, European Liquidity Centre operating policies~; LSEG, Data Centre Migration
plan (2022).
\item McKay Brothers, communiqué du 10 mars 2015~; codes d'identification de marché ISO 10383~; coordonnées des sites issues d'OpenStreetMap Nominatim.
\end{omsources}

\section{Exercices}

\begin{exercise}[$\star$]\label{exo:nw:the-european-map:1}
Calculer de combien a changé le plancher à sens unique en fibre, depuis Slough LD4 jusqu'au centre de données principal d'Euronext, le 6 juin 2022.
\end{exercise}

\begin{exercise}[$\star$]\label{exo:nw:the-european-map:2}
Le chiffre 2015 de McKay Brothers pour LD4--FR2 était un aller-retour de \qty{4.64}{\milli\second}. Calculer son \omterm{def:nw:the-north-american-map:factor}{facteur de tracé} dans l'air.
\end{exercise}

\begin{exercise}[$\star$]\label{exo:nw:the-european-map:3}
Expliquer la différence entre un \omterm{def:nw:the-european-map:pop}{point de présence} et la \omterm{def:nw:colocation-products-and-how-they-are-sold:colo}{colocation}.
\end{exercise}

\begin{exercise}[$\star\star$]\label{exo:nw:the-european-map:4}
Pourquoi la FAQ d'Euronext promettait-elle une égalisation «~depuis l'arrière de l'antenne jusqu'à la Meet Me Room~»~? Que se passerait-il sans elle~?
\end{exercise}

\begin{exercise}[$\star\star$]\label{exo:nw:the-european-map:5}
Une société relaie des signaux entre Eurex et Euronext. Comparer le chemin passant par Slough avec celui passant par Zurich, et avec le chemin direct.
\end{exercise}

\begin{exercise}[$\star\star$]\label{exo:nw:the-european-map:6}
EWIN est publié comme «~moins de \qty{4}{\milli\second}~». Que cela vous apprend-il, et que ne vous apprend-il pas~?
\end{exercise}

\begin{exercise}[$\star\star\star$]\label{exo:nw:the-european-map:7}
\emph{Programmation.} Ajouter une ligne au registre qui donne à Euronext Paris deux sites principaux le même jour, puis demander son site principal. Que se passe-t-il,
et pourquoi est-ce le bon comportement~?
\end{exercise}

\begin{exercise}[$\star\star\star$]\label{exo:nw:the-european-map:8}
\emph{Trouver la faille.} «~Depuis la migration, les sociétés londoniennes accusent trois millisecondes de retard sur les sociétés italiennes sur Euronext, si bien que les teneurs de marché londoniens ne peuvent pas y rivaliser.~»

\end{exercise}

\section{Problème~: de Basildon à Bergame}

\begin{problem}[{Problème du week-end --- ce que la migration d'Euronext a fait à la carte}]\label{pb:nw:the-european-map:1}
Un teneur de marché cote les marchés d'Euronext depuis Slough LD4 et se couvre sur Eurex à Francfort FR2. Le 6 juin 2022, le centre de données principal d'Euronext
passe de Basildon à Bergame.

\medskip\noindent\textbf{Partie I --- Avant.}
\begin{enumerate}
\item Quelles étaient la \omterm{def:nw:the-north-american-map:geodesic}{distance géodésique} et les planchers à sens unique (vide, fibre) de LD4 à Basildon~?
\item Et de Basildon à FR2~?
\item Quelle était la forme du triangle LD4--FR2--Euronext~?
\item Où la société pourrait-elle placer un serveur pour être proche à la fois d'Euronext et des autres plateformes londoniennes~?
\end{enumerate}

\medskip\noindent\textbf{Partie II --- Après.}
\begin{enumerate}[resume]
\item Quelles sont les distances et les planchers de LD4 à Bergame, dans le vide et dans la fibre~?
\item De combien chaque plancher a-t-il changé~?
\item De FR2 à Bergame, et de combien cela a-t-il changé~?
\item Quelle est la nouvelle forme du triangle~?
\end{enumerate}

\medskip\noindent\textbf{Partie III --- Les choix de la société.}
\begin{enumerate}[resume]
\item La société garde ses serveurs à LD4 et rejoint Bergame par fibre avec un \omterm{def:nw:the-north-american-map:factor}{facteur de tracé} de 1.2. Quelle est sa latence à sens unique~?
\item Elle achète EWIN à la place. Au plus, combien de temps prennent ses ordres, et que gagne-t-elle par rapport à la réponse de la première question de la partie III~?
\item Son moteur de cotation pour les marchés d'Euronext doit-il passer à Bergame~? Que perd-il si c'est le cas~?
\item À quelle distance de la géodésique Francfort--Bergame se trouve Zurich, en détour~?
\end{enumerate}

\medskip\noindent\textbf{Partie IV --- Le verdict.}
\begin{enumerate}[resume]
\item Énoncer le \emph{résultat nommé}~: la variation du plancher à sens unique de LD4 aux moteurs d'appariement d'Euronext, et le nouveau triangle.
\item Pourquoi Euronext a-t-elle ouvert des \omterm{def:nw:the-european-map:pop}{points de présence} à Londres avant la migration~?
\item Que consigne le registre, et pourquoi deux lignes plutôt qu'une ligne modifiée~?
\item Quel autre déménagement européen de 2022--2023 le registre devrait-il consigner~?
\item Que dit la \cref{prop:nw:the-european-map:triangle} sur l'endroit où placer un relais entre Eurex et Euronext~?
\item Quels sites du chapitre ne sont qu'au niveau de la rue, et cela importe-t-il pour le résultat nommé~?
\item Que changerait une fibre à cœur creux (chapitre~13) sur le couloir Londres--Bergame~?
\item En une phrase~: que la migration a-t-elle changé pour une société londonienne~?
\end{enumerate}
\end{problem}

\section{Questions d'entretien}

\begin{interviewq}[$\star$ \iqroles{développeur, trader}]\label{iq:nw:the-european-map:1}
Où se trouvent les principaux moteurs d'appariement européens, et quel est le \omterm{def:nw:the-north-american-map:floor}{plancher de latence} Londres--Francfort~?
\end{interviewq}

\begin{interviewq}[$\star$ \iqroles{développeur}]\label{iq:nw:the-european-map:2}
Qu'est-ce qu'un \omterm{def:nw:the-european-map:pop}{point de présence}, et pourquoi une bourse en ouvrirait-elle un dans un autre pays~?
\end{interviewq}

\begin{interviewq}[$\star\star$ \iqroles{développeur}]\label{iq:nw:the-european-map:3}
Comment tiendriez-vous à jour une table indiquant quelle plateforme se trouve dans quel centre de données~?
\end{interviewq}

\begin{interviewq}[$\star\star$ \iqroles{trader, développeur}]\label{iq:nw:the-european-map:4}
Une bourse annonce qu'elle va déplacer son moteur d'appariement dans un autre pays. Que faites-vous, et dans quel ordre~?
\end{interviewq}

\begin{interviewq}[$\star\star$ \iqroles{trader}]\label{iq:nw:the-european-map:5}
Pourquoi les bourses vendent-elles des liaisons hertziennes micro-ondes à leurs propres centres de données~?
\end{interviewq}

\begin{interviewq}[$\star\star\star$ \iqroles{développeur, chercheur}]\label{iq:nw:the-european-map:6}
Où placeriez-vous des serveurs pour trader Eurex contre Euronext, et comment vérifieriez-vous ce choix~?
\end{interviewq}
